% Eine Quelle, zwei PDFs. build.sh setzt \explainmode:
%   0 -> loesung_blatt01.pdf
%   1 -> erklaerung_blatt01.pdf
\documentclass[a4paper,11pt]{article}
\usepackage[T1]{fontenc}
\usepackage[utf8]{inputenc}
\usepackage[english]{babel}
\usepackage{amsmath,amssymb,mathtools}
\usepackage[margin=2.2cm]{geometry}
\usepackage{enumitem}
\usepackage{xcolor}
\usepackage{booktabs}
\usepackage{tikz}
\usepackage{pgfplots}\pgfplotsset{compat=1.18}
\usepackage[most]{tcolorbox}
\setlength{\parindent}{0pt}
\setlength{\parskip}{4pt}

\providecommand{\explainmode}{0}
\newif\ifexplain
\ifnum\explainmode=1 \explaintrue\fi

\newcommand{\blatttitel}[2]{\begin{center}{\Large\bfseries #1}\\[2pt]#2\end{center}}
\newcommand{\aufgabe}[1]{\section*{#1}}
\newcommand{\angabe}[1]{\textit{#1}\par\smallskip}
\newcommand{\teil}[1]{\subsection*{(#1)}}
\newcommand{\why}[1]{\text{\small\color{black!65}#1}}

\newtcolorbox{erklbox}{colback=blue!4,colframe=blue!40,breakable,boxrule=0.5pt,
  fonttitle=\bfseries\small,title={Erklärung}}
\newcommand{\explain}[1]{\ifexplain\begin{erklbox}#1\end{erklbox}\fi}
\newcommand{\explainfig}[1]{\ifexplain\begin{center}#1\end{center}\fi}

\begin{document}
\blatttitel{Exercise sheet 1 -- \ifexplain Explanation\else Solution\fi}{Linear Optimization (WS 26/27)}

\explain{Dieses Blatt hat drei Typen von Aufgaben: Bei der Matrixaufgabe wird Invertierbarkeit begründet. Bei den Begriffen wird eine KI-Antwort mit den Vorlesungsdefinitionen verglichen. Bei den letzten zwei Aufgaben wird jeweils ein Optimierungsmodell formuliert: zuerst Variablen, dann Zielfunktion, dann Nebenbedingungen.}

\aufgabe{Classroom task 1.1}
\angabe{Consider the matrix
\[
A=\begin{pmatrix}
0&1&0\\
2&0&1\\
1&1&0
\end{pmatrix}.
\]
How many ways do you know to characterize invertibility of $A$? Which one is the most useful to see or prove invertibility of $A$?}

\explain{Für eine quadratische Matrix gibt es mehrere gleichwertige Kennzeichen für Invertierbarkeit. Am schnellsten ist hier die Determinante, weil eine $3\times3$-Determinante direkt ausgerechnet werden kann. Wenn die Determinante nicht $0$ ist, ist die Matrix invertierbar.}

Some equivalent characterizations for a square $3\times3$ matrix are:
\begin{itemize}[leftmargin=2em]
  \item $A$ has an inverse matrix $A^{-1}$.
  \item $\det(A)\ne 0$.
  \item $\operatorname{rank}(A)=3$.
  \item The rows of $A$ are linearly independent.
  \item The columns of $A$ are linearly independent.
  \item The homogeneous system $Ax=0$ has only the solution $x=0$.
  \item For every $b\in\mathbb{R}^3$, the system $Ax=b$ has exactly one solution.
  \item The columns of $A$ span $\mathbb{R}^3$.
  \item $0$ is not an eigenvalue of $A$.
\end{itemize}

For the given matrix, the determinant test is useful:
\begin{align*}
\det(A) &= 0\cdot0\cdot0 + 1\cdot1\cdot1 + 0\cdot2\cdot1 - 0\cdot0\cdot1 - 1\cdot2\cdot0 - 0\cdot1\cdot1 && \why{Sarrus rule}\\
&= 0 + 1 + 0 - 0 - 0 - 0 && \why{multiply all factors}\\
&= 1 && \why{add and subtract the terms}
\end{align*}

Check by expansion along the first row:
\begin{align*}
\det(A) &= -\det\begin{pmatrix}2&1\\1&0\end{pmatrix} && \why{only the entry $a_{12}=1$ contributes}\\
&= -\bigl(2\cdot0-1\cdot1\bigr) && \why{determinant of a $2\times2$ matrix}\\
&= -(0-1) && \why{$2\cdot0=0$ and $1\cdot1=1$}\\
&= 1 && \why{$0-1=-1$ and $-(-1)=1$}
\end{align*}

\[
\boxed{\det(A)=1\ne0,\text{ so }A\text{ is invertible.}}
\]

\aufgabe{Exercise 1.1}
\angabe{(1) Ask an AI of your choice the following question: ``What is a linear program in standard form?'' Compare the answer to the definition of standard form in the lecture notes. (2) Continue by asking the AI about each of the definitions of feasible or infeasible vector, feasible set, feasible or infeasible linear program, optimal value, bounded or unbounded linear program and optimal solution. Do the answers match the definitions in the lecture notes? For the submission, copy or write down the key parts of the answers of the AI and just tick it, if it is the same definition and if not, comment on the differences.}

\explain{Die Vorlesung verwendet für die Standardform eine Maximierungsaufgabe mit Gleichungsnebenbedingungen und Nichtnegativitätsbedingungen. Viele KI-Antworten nennen zusätzlich Minimierung oder Ungleichungen. Das ist nicht falsch als allgemeine Beschreibung, aber es ist nicht exakt dieselbe Normalisierung wie in den Vorlesungsnotizen.}

\teil{1}
Key part of the AI answer:
\[
\begin{array}{ll}
\text{maximize} & c^T x\\
\text{subject to} & Ax=b,\\
& x\ge 0,
\end{array}
\]
where $x\in\mathbb{R}^n$ is the decision variable, $c\in\mathbb{R}^n$, $A\in\mathbb{R}^{m\times n}$ and $b\in\mathbb{R}^m$ are given data.

Comparison with the lecture notes:
\[
\boxed{\text{This matches the lecture definition of standard form. }\checkmark}
\]

Comment: If an AI answer also allows minimization problems or inequality constraints in the displayed standard form, then it is more general than the lecture definition. In the lecture notes, these forms are transformed into the standard form above by changing signs, adding slack variables, and replacing free variables by the difference of two nonnegative variables.

\explain{Wichtig ist der Unterschied zwischen einer allgemeinen linearen Optimierungsaufgabe und der Standardform der Vorlesung. Die Standardform ist eine feste Schreibweise. Andere Schreibweisen können äquivalent sein, müssen aber erst umgeformt werden.}

\teil{2}
The key definitions from the AI answer and the comparison with the lecture notes are:

\begin{center}
\begin{tabular}{p{0.26\linewidth}p{0.56\linewidth}p{0.10\linewidth}}
\toprule
Term & Key part of the AI answer & Same as lecture notes?\\
\midrule
feasible vector & A vector $x\in\mathbb{R}^n$ is feasible if it satisfies all constraints $Ax=b$ and $x\ge0$; otherwise it is infeasible. & $\checkmark$\\
feasible set & The feasible set is $C=\{x\in\mathbb{R}^n\mid Ax=b,\ x\ge0\}$. & $\checkmark$\\
feasible linear program & A linear program is feasible if $C\ne\emptyset$. & $\checkmark$\\
infeasible linear program & A linear program is infeasible if $C=\emptyset$. & $\checkmark$\\
optimal value & The optimal value is $v=\sup\{c^T x\mid x\in C\}$; for an infeasible LP the convention is $v=-\infty$. & $\checkmark$\\
unbounded linear program & A feasible LP is unbounded if for every $M\in\mathbb{R}$ there is an $x\in C$ with $c^Tx>M$; then $v=+\infty$. & $\checkmark$\\
bounded linear program & A feasible LP is bounded if there is an $M\in\mathbb{R}$ such that $c^Tx\le M$ for every $x\in C$. & $\checkmark$\\
optimal solution & A feasible vector $x^\ast\in C$ is an optimal solution if $c^Tx^\ast\ge c^Tx$ for every $x\in C$. & $\checkmark$\\
\bottomrule
\end{tabular}
\end{center}

\[
\boxed{\text{All listed definitions match the lecture notes.}}
\]

\explain{Bei Maximierungsproblemen wird nach oben verglichen: Eine optimale Lösung hat mindestens so großen Zielfunktionswert wie jede andere zulässige Lösung. Deshalb steht bei der optimalen Lösung $\ge$ und bei Beschränktheit nach oben $\le M$.}

\aufgabe{Exercise 1.2}
\angabe{A student has a total of 30 minutes to figure out the course requirements. The sources have the following values and times: asking a colleague: value $4$, time $5$ min; coming to the first lecture: value $8$, time $9$ min; reading the description in the Moodle course: value $7$, time $7$ min; reading the LV-Karte: value $16$, time $15$ min; doing nothing to not risk the headache: value $40$, time $30$ min. Since we do not want to rely only on one source, we need to use at least two. The goal is to maximize the personal value within the time constraints. Formulate the optimization problem.}

\explain{Jede Informationsquelle wird entweder verwendet oder nicht verwendet. Daher passen Binärvariablen: $1$ bedeutet „verwenden“, $0$ bedeutet „nicht verwenden“. Die Nebenbedingung mit der Summe der Variablen erzwingt mindestens zwei Quellen.}

Let
\[
 x_i=\begin{cases}
 1, & \text{if source }i\text{ is used},\\
 0, & \text{if source }i\text{ is not used},
 \end{cases}
 \qquad i=1,\ldots,5.
\]
The indices correspond to the five sources in the order stated in the exercise.

The objective coefficients are the values:
\[
(4,8,7,16,40).
\]
The time coefficients are the required minutes:
\[
(5,9,7,15,30).
\]
The time limit is $30$ minutes and at least two sources must be used. Hence the model is
\[
\boxed{
\begin{array}{ll}
\text{maximize} & 4x_1+8x_2+7x_3+16x_4+40x_5\\[2pt]
\text{subject to} & 5x_1+9x_2+7x_3+15x_4+30x_5\le 30,\\
& x_1+x_2+x_3+x_4+x_5\ge 2,\\
& x_i\in\{0,1\},\qquad i=1,\ldots,5.
\end{array}}
\]

Check of the modelling choices:
\begin{itemize}[leftmargin=2em]
  \item The first constraint counts minutes and compares them with $30$.
  \item The second constraint counts the selected sources and requires at least $2$.
  \item The binary condition means that every source is either selected once or not selected.
\end{itemize}

\explain{Die fünfte Möglichkeit hat schon allein $30$ Minuten. Zusammen mit einer zweiten Quelle würde sie die Zeitgrenze überschreiten. Man muss sie trotzdem nicht vorher streichen: Das Modell macht diese Kombinationen automatisch unzulässig.}

\[
\boxed{\text{The optimization problem is the binary linear program displayed above.}}
\]

\aufgabe{Exercise 1.3}
\angabe{The company ``FreshShine'' produces three products. Product 1, Citrus Shower Gel, yields a profit of $3$ euros per unit. Product 2, Citrus Shampoo, yields a profit of $5$ euros per unit. Product 3, Citrus Body Lotion, yields a profit of $4$ euros per unit. Production requires machines A, B and C. Shower Gel needs $1$ hour on A and $1$ hour on B. Shampoo needs $2$ hours on A, $1$ hour on B and $3$ hours on C. Body Lotion needs $1$ hour on A, $2$ hours on B and $1$ hour on C. The available capacities are $170$ hours on A, $150$ hours on B and $180$ hours on C. The company can sell all units produced. The goal is to maximize profit. Formulate the mathematical optimization model.}

\explain{Hier sind die Variablen Produktionsmengen. Die Zielfunktion ist der Gewinn. Jede Maschine liefert eine eigene Kapazitätsnebenbedingung: Verbrauch pro Stück mal Stückzahl, aufsummiert über alle Produkte, darf die verfügbare Zeit nicht überschreiten.}

Let
\[
\begin{array}{ll}
x_1 & = \text{number of units of Citrus Shower Gel produced per month},\\
x_2 & = \text{number of units of Citrus Shampoo produced per month},\\
x_3 & = \text{number of units of Citrus Body Lotion produced per month}.
\end{array}
\]

Profit per month is
\[
3x_1+5x_2+4x_3.
\]
Machine A is used for
\[
1x_1+2x_2+1x_3
\]
hours. Machine B is used for
\[
1x_1+1x_2+2x_3
\]
hours. Machine C is used for
\[
0x_1+3x_2+1x_3
\]
hours.

Therefore the linear optimization model is
\[
\boxed{
\begin{array}{ll}
\text{maximize} & 3x_1+5x_2+4x_3\\[2pt]
\text{subject to} & x_1+2x_2+x_3\le170,\\
& x_1+x_2+2x_3\le150,\\
& 3x_2+x_3\le180,\\
& x_1,x_2,x_3\ge0.
\end{array}}
\]

Check of the modelling choices:
\begin{itemize}[leftmargin=2em]
  \item The coefficient row $(1,2,1)$ is the monthly use of machine A per unit of products $1,2,3$.
  \item The coefficient row $(1,1,2)$ is the monthly use of machine B per unit of products $1,2,3$.
  \item The coefficient row $(0,3,1)$ is the monthly use of machine C per unit of products $1,2,3$.
  \item The nonnegativity constraints exclude negative production quantities.
\end{itemize}

\explain{Die Aufgabe verlangt nur das Modell. Wenn ganze Stückzahlen vorgeschrieben wären, müsste zusätzlich $x_1,x_2,x_3\in\mathbb{Z}$ verlangt werden. Da das Blatt von linearer Optimierung spricht und keine Ganzzahligkeit fordert, wird das kontinuierliche lineare Modell verwendet.}

\[
\boxed{\text{The mathematical optimization model is the LP displayed above.}}
\]

\end{document}
